% !TEX program = xelatex
% 編譯：xelatex final_report.tex（需 ctex / xeCJK 環境，建議 TeX Live 2022+）
\documentclass[11pt, fontset=none]{ctexart}
% fontset=none：改用系統 Noto CJK TC，避免 fandol 缺繁體字形（豆腐字）
\setCJKmainfont[AutoFakeSlant=0.2]{Noto Serif CJK TC}
\setCJKsansfont{Noto Sans CJK TC}
\setCJKmonofont{Noto Sans Mono CJK TC}

\usepackage{geometry}
\geometry{a4paper, margin=2.5cm}
\usepackage{amsmath, amssymb}
\usepackage{booktabs}
\usepackage{graphicx}
\usepackage{xcolor}
\usepackage{listings}
\usepackage{hyperref}
\hypersetup{colorlinks=true, linkcolor=blue!55!black, urlcolor=blue!55!black,
            citecolor=blue!55!black}
\usepackage{enumitem}
\usepackage{tikz}
\usetikzlibrary{arrows.meta, positioning, shapes.geometric}

\lstset{
  basicstyle=\ttfamily\footnotesize,
  breaklines=true,
  frame=single,
  framesep=4pt,
  rulecolor=\color{black!25},
  backgroundcolor=\color{black!3},
  columns=fullflexible,
}

\title{\textbf{深度學習投資組合自動交易機器人\\—— 台股 88 檔宇宙之 LSTM 集成策略設計與實證}}
\author{自動交易專題 期末成果報告}
\date{2026 年 6 月}

\begin{document}
\maketitle

\begin{abstract}
  \noindent
  本專題設計一套\textbf{每日自動再平衡}的台股投資組合交易機器人。核心為一個\emph{跨個股共享、置換等變}的 LSTM 評分網路（\texttt{PortfolioNetLSTM}），以 50 日視窗、每檔 14 維特徵（8 維技術指標 + 4 維橫斷面排名 + 波動／動能）對每檔股票輸出分數，經 softmax 與\textbf{每檔上限（per-stock cap）}得到目標權重，並以\textbf{負 Sharpe（含換手成本）}端到端訓練。模型歷經三代演進：v4（中位數集成）$\rightarrow$ v5（Sharpe 加權集成 + 注水式溢出重分配 water-fill + 集成 Top-$N$ 截斷）$\rightarrow$ v6（股票宇宙由 46 檔擴充至 88 檔、Top-$N=10$ 強化集中度、每日重訓）。在 2019–2023 共 10 折\emph{滾動式前向驗證（walk-forward）}下，v5 相對等權基準之 OOS 超額報酬中位數為 \textbf{+2.71\%}（平均 +6.96\%），勝過 v4 的 +0.98\%；v6 在可驗證的 7 折上中位超額報酬達 \textbf{+5.65\%}。關鍵發現為：\emph{更大的候選池需搭配更強的集中度}（Top-10 而非 Top-15）方能濾除雜訊。系統已於 GitHub Actions 上以每日 17:00（台北）排程實盤（模擬）部署，並內建回撤停損與最小交易次數（$\geq 100$ 筆）之湊單機制。
\end{abstract}

\section{專題目標與條件}

\subsection{給定條件}
\begin{itemize}[leftmargin=2em]
  \item 初始資本：100{,}000{,}000 元（虛擬資金）。
  \item 資料來源：近十年每日台股歷史交易資訊（漲跌幅、開、收、高、低、成交量）。
  \item 交易介面：股票買賣交易 API（NCKU \texttt{sim\_stock}）。
\end{itemize}

\subsection{目標}
設計可\textbf{每日執行}的自動化交易機器人，於可交易期間（即日起至 2026/06/26）內最大化投資利潤，且系統須\textbf{成功自動完成至少 100 筆交易}。投資利潤定義為
\[
  \text{利潤} \;=\; \big[\text{2026/06/26 收盤後資產總額}\big] \;-\; 100{,}000{,}000.
\]

\section{系統架構概覽}

整體為三個彼此獨立、僅共用交易 API 的子系統：\textbf{資料層}（\texttt{stock\_api/}，直抓 TWSE／TPEX 官方價量，統一 10 欄 schema，內含指數退避、ROC 日期轉換與除權息清理）、\textbf{模型／部署層}（\texttt{RL/}，特徵工程 $\to$ LSTM 評分 $\to$ 集成決策 $\to$ 實盤下單）、與 \textbf{CI 自動化}（兩條 GitHub Actions：工作日 17:00 台北每日部署 + 每日重訓，狀態與權重提交回 \texttt{main}）。每日交易管線如圖~\ref{fig:pipeline}；技術細節集中於後續章節，此處僅勾勒資料流。

\begin{figure}[h]
  \centering
  \begin{tikzpicture}[
      node distance=8mm and 12mm,
      box/.style={draw, rounded corners, align=center, minimum height=8mm,
          minimum width=22mm, font=\small, fill=black!4},
      >=Stealth]
    \node[box] (twse) {TWSE / TPEX\\ 官方價量};
    \node[box, right=of twse] (csv) {本地快取\\ \texttt{RL/data/*.csv}};
    \node[box, right=of csv] (feat) {特徵工程\\ 14 維 / 檔};
    \node[box, right=of feat] (lstm) {LSTM 評分\\ + 集成};
    \node[box, below=of lstm] (order) {權重 $\to$ 整股\\ 下單（含費稅）};
    \node[box, left=of order] (broker) {交易 API\\ 對帳 / 成交};
    \node[box, left=of broker] (state) {帳本狀態\\ \texttt{deploy\_state.json}};
    \draw[->] (twse)--(csv); \draw[->] (csv)--(feat); \draw[->] (feat)--(lstm);
    \draw[->] (lstm)--(order); \draw[->] (order)--(broker); \draw[->] (broker)--(state);
    \draw[->] (state.north) |- ([yshift=4mm]feat.north) -| (feat.north);
  \end{tikzpicture}
  \caption{每日交易管線：抓取／快取 $\to$ 特徵 $\to$ LSTM 集成評分 $\to$ 整股下單 $\to$ 對帳 $\to$ 帳本持久化。}
  \label{fig:pipeline}
\end{figure}

\section{模型與核心技術}

\subsection{PortfolioNetLSTM}
網路對每檔股票之 50 日特徵窗（含 4 維可學習個股嵌入 \texttt{nn.Embedding}）跑\emph{共享權重}的單層 LSTM，取末步隱狀態經線性層輸出一個分數；$N$ 檔分數與一個可學習現金分數串接後 softmax，得到含現金槽的權重向量，再施以每檔上限 \texttt{max\_weight}$=0.10$。LSTM 權重跨個股共享、與 $N$ 幾乎無關，僅嵌入表 $N\times 4$ 隨宇宙線性微增，故由 46 擴至 88 在計算上乾淨可行——此\emph{置換等變}設計正是 v6 能無痛擴池的前提。

\subsection{特徵工程}
每檔股票每日構造 14 維特徵：8 維\emph{個股本地技術指標}（\texttt{return}、\texttt{bias\_5}、\texttt{bias\_20}、\texttt{macd\_h}、\texttt{rsi\_14}、\texttt{bb\_pos}、\texttt{atr}、\texttt{capacity\_change}）、4 維\emph{橫斷面每日排名}（return、rsi、bias\_20、量能變化之 pct-rank），以及 v6 新增之 \texttt{vol\_20}（20 日對數報酬標準差）與 \texttt{mom\_60}（60 日動能）。橫斷面排名提供「市場相對」訊號，是模型超額報酬的關鍵來源。

\subsection{股票宇宙}
v4／v5 使用 46 檔台股大型權值股；v6 依台灣 50（0050）與中型 100（0051／0100）成分股，篩除興櫃（ESB，無真實開盤價）後擴充至 \textbf{88 檔}（原 46 + 43 檔流動性佳之上市／上櫃中大型股；另剔除 2024 才上市、歷史過短的 6446）。個股順序為\emph{承載性的（load-bearing）}：檢查點儲存訓練時的順序，部署端比對不符即拒絕執行，故擴充宇宙必須完整重訓。

\subsection{損失函數：成本感知的負 Sharpe}
端到端最小化\textbf{負 Sharpe}，並在報酬序列中計入 L1 換手成本：
\[
  \text{pnl}_t \;=\; \langle w_t,\, r_t\rangle \;-\; c\cdot \big\| w_t - \tilde{w}_{t-1}\big\|_1,
\]
其中 $\tilde{w}_{t-1}$ 為前期權重經當期資產漂移後之值，$c=0.0042$。模型於訓練時即「感知」交易成本，自發傾向降低不必要換手，無須額外正則項。

\section{開發歷程：三代演進與關鍵槓桿}

模型並非一次成形，而是\emph{由實證驅動}的三代迭代：每一代都源自前代在 walk-forward 上暴露的具體缺陷，再以一個針對性的技術槓桿修正（表~\ref{tab:evo}）。

\begin{table}[h]
  \centering
  \caption{三代演進路線圖：各代的關鍵技術改動、動機與成效。}
  \label{tab:evo}
  \small
  \begin{tabular}{lll c}
    \toprule
    代  & 關鍵技術改動                          & 動機／解決的問題          & OOS 中位超額 \\
    \midrule
    v4 & 5-seed 逐維中位數集成                  & 降低單模型變異          & 基線（抹除高信念部位） \\
    v5 & 注水溢出 + Sharpe 加權集成 + Top-15 截斷 & 消除現金拖累、恢復集中度     & +2.71\%       \\
    v6 & 宇宙 46$\to$88 + Top-10 + 每日重訓    & 擴大 alpha 源、濾除雜訊  & \textbf{+5.65\%}（7 折） \\
    \bottomrule
  \end{tabular}
\end{table}

\subsection{v4：基線與「中位數抹除」問題}
5 seed 以逐維中位數集成以求穩健，診斷卻發現中位數會把任一 seed 為零的部位抹平——而那正是模型最高信念的重押：單 seed walk-forward 超額達 +4.09\%，中位集成卻塌陷為 +0.98\%。\emph{集成方式本身吃掉了 alpha}，成為下一代的首要攻克對象。

\subsection{v5 槓桿 A：注水式溢出重分配（water-fill）}
問題：超過每檔上限的權重原本一律倒入現金，使多頭期間約 48\% 資金閒置、結構性落後等權。逐 fold 拆解坐實病灶——fold-0（強多頭）現金高達 54.7\%。解法不是取消上限（會回到過度集中），而是把\emph{機械式溢出}按現價權重重分配給未達上限的個股，同時\emph{保留}模型刻意持有的現金。效果立竿見影：fold-0 現金 54.7\%$\to$10.5\%、整體現金 $\sim$48\%$\to\sim$3\%，超額升至 +4.27\%（平均，10 折），證實現金外溢確為當時主要拖累。

\subsection{v5 槓桿 B：Sharpe 加權集成 + 集成 Top-$N$ 截斷}
以 $\mathrm{softmax}(\text{val Sharpe})$ 對 seed 加權取代中位數，讓驗證表現好的 seed 主導；再於\emph{集成後}硬性保留權重最大的 $N$ 檔、其餘歸零重標準化，找回「敢重押」的集中度。工程上，三處上限／截斷邏輯（訓練前傳、walk-forward 集成、部署集成）以單一共享函數鎖定，並以對拍測試保證 \emph{backtest = live}，杜絕「回測賺、實盤賠」的口徑漂移。v5 最終 OOS 中位超額 +2.71\%（平均 +6.96\%）、7/10 折正報酬，全面超越 v4。

\subsection{v6：更大的池、更少的部位}
假設更大的選股池能提供更多 alpha 源，遂將宇宙由 46 擴至 88 檔並改為每日重訓。但首次驗證\emph{反而退步}（Top-15 中位 $-0.01\%$）：更大的池同時引入更多雜訊與邊際個股。修正來自 Top-$N$ 離線掃描（見 \S\ref{sec:topn}）——須\textbf{收緊至 Top-10} 取最高信念的 10 檔，中位超額才躍升至 +5.65\%。核心教訓：\emph{更大的候選池必須搭配更強的集中度}。

\section{交易策略與部署}

\begin{itemize}[leftmargin=2em]
  \item \textbf{每日再平衡}：以 5 seed 集成輸出之目標權重，與券商對帳後之實際庫存比較，差額轉為\textbf{整股（1000 股／張）}買賣單；買單以 $\min(\text{開},\text{收})$、賣單以 $\max(\text{開},\text{收})$ 報價以提高成交率，價格依台股 tick 規則四捨五入，計入 0.1425\% 手續費與 0.3\% 證交稅（賣方），最低 20 元。
  \item \textbf{回撤停損}：以峰值資產為基準，若 $(\text{peak}-\text{port})/\text{peak} > 10\%$ 即全數平倉並\emph{黏滯停機（sticky halt）}，需人工解除——此為實盤的最終風險防線。
  \item \textbf{對帳優先}：\texttt{--live} 每次先呼叫 \texttt{Get\_User\_Stocks} 以券商真實庫存覆寫本地庫存，再行決策（「先知道手上有多少股與現金，再下單」），可校正被拒單造成的漂移。
  \item \textbf{最小交易次數（湊單）}：為滿足「$\geq 100$ 筆交易」之系統要求，於模型再平衡後追加 \texttt{--pad-pairs}$=12$ 之\emph{賣出—買回}往返單（每對 2 筆、對組合淨中性，僅付單邊約 0.585\% 費稅於一張最便宜持股），每個交易日穩定新增 24 筆，數個交易日內即可達標。
\end{itemize}

\section{實驗流程}

\subsection{滾動式前向驗證（Walk-forward）}
採 2019 H1 至 2023 H2 共 \textbf{10 個 6 個月折}，每折以擴張訓練集（自 2015 起）之最近 500 個交易日訓練 5 seed，再於該折驗證窗計算超額報酬、Sharpe、現金佔比與活躍檔數。集成決策規則與實盤完全相同（同一 \texttt{aggregate\_weights}），故 OOS 數字與部署決策「同口徑」。

\subsection{Top-$N$ 離線掃描（探針）}\label{sec:topn}
由於 Top-$N$ 截斷為\emph{集成後}之後處理、不改變訓練，故可將一次 walk-forward 快取之各 seed 每日權重\textbf{離線重放}於不同 $N$（\texttt{probe\_ensemble\_topn.py}），於數秒內掃出最佳 $N$，無須重訓。此設計使「更多選擇」這個假設能以近乎零成本驗證。

\section{結果分析}

\subsection{各代 OOS 超額報酬（vs.\ 等權基準）}
表~\ref{tab:gen} 比較三代模型之 walk-forward 表現。v5 的注水 + Top-15 顯著優於 v4；v6 擴充宇宙後，\emph{在 Top-15 反而退步}，須\textbf{加強集中度至 Top-10} 方勝出。

\begin{table}[h]
  \centering
  \caption{各代模型之 walk-forward OOS 超額報酬（相對等權基準）。}
  \label{tab:gen}
  \small
  \begin{tabular}{lccccc}
    \toprule
    模型         & 宇宙 & 集成／Top-$N$           & 平均超額             & 中位超額             & 正報酬折數 \\
    \midrule
    v4         & 46 & 中位數                  & +0.98\%          & ---              & 5/10  \\
    v4（單 seed） & 46 & ---                  & +4.09\%          & ---              & 5/10  \\
    v5         & 46 & 注水 + Top-15          & \textbf{+6.96\%} & +2.71\%          & 7/10  \\
    v6（7 折）    & 88 & 注水 + Top-15          & +3.84\%          & $-0.01\%$        & 3/7   \\
    v6（7 折）    & 88 & 注水 + \textbf{Top-10} & +6.11\%          & \textbf{+5.65\%} & 4/7   \\
    \bottomrule
  \end{tabular}
\end{table}

\subsection{核心發現：更大的池、更少的部位}
表~\ref{tab:topn} 為 v6 快取上之 Top-$N$ 掃描。當候選池由 46 擴至 88，雜訊增加，\emph{放寬}選股（$N\geq15$，活躍 15--34 檔）反而拖入邊際個股而表現平庸；\emph{收緊}至 Top-10 取最高信念之 10 檔，中位超額由 $-0.01\%$ 躍升至 \textbf{+5.65\%}。但過度集中（$N\leq8$）又使現金外溢、中位塌陷。最佳點落在 $N\approx10$。

\begin{table}[h]
  \centering
  \caption{v6（88 檔）集成 Top-$N$ 離線掃描（7 折）。}
  \label{tab:topn}
  \small
  \begin{tabular}{cccccc}
    \toprule
    Top-$N$     & 平均超額             & 中位超額             & 正報酬折數 & 平均現金   & 平均活躍檔數 \\
    \midrule
    none(34)    & $-0.06\%$        & $-1.42\%$        & 3/7   & 3.4\%  & 34.2   \\
    20          & +2.52\%          & $-1.72\%$        & 3/7   & 4.2\%  & 20.0   \\
    15          & +3.84\%          & $-0.01\%$        & 3/7   & 4.8\%  & 15.0   \\
    12          & +5.22\%          & +2.97\%          & 4/7   & 5.4\%  & 12.0   \\
    \textbf{10} & \textbf{+6.11\%} & \textbf{+5.65\%} & 4/7   & 5.9\%  & 10.0   \\
    8           & +5.14\%          & +3.46\%          & 4/7   & 21.9\% & 8.0    \\
    \bottomrule
  \end{tabular}
\end{table}

\subsection{實盤部署}
v6 已於 2026/06/18 完成上線（\texttt{main}），首個交易日依 Top-10 規則建立 10 檔持股、現金近乎全數投入、回撤 0\%（峰值重設，不繼承前代回撤）。每日由 CI 自動再平衡並追加湊單，最終損益以 2026/06/26 收盤結算。

\section{討論與限制}
\begin{itemize}[leftmargin=2em]
  \item \textbf{資料對齊的擴展性瓶頸}：\texttt{build\_tensors} 以各股日期\emph{交集}對齊，在 88 檔下各股零星缺值\emph{複合}累積，使近期 3 折（2022--2023）共同日數塌陷至無法成窗，故 v6 僅在 0--6 折（7 折）完成驗證。根因之正解為改採\emph{聯集 + 前向填補 + 遮罩}，可解鎖全 10 折驗證並支援宇宙進一步成長——列為首要後續工作。（注：部署端推論為逐股、含遮罩，故不受此影響。）
  \item \textbf{高變異與 fold 依賴}：超額報酬標準差約 24\%，平均值常受單一強勢 fold 拉抬；中位數方為可信指標，且 v6 僅 7 折驗證，需以上線初期實盤行為交叉佐證。
  \item \textbf{帳戶起點}：v6 繼承前代 v4 觸發停損後之資本（約 74.8M，相對 100M 名目已 $-25\%$）；本報告之模型優勢以\emph{OOS 超額報酬}論證，最終競賽損益則受此起點與短窗影響。
\end{itemize}

\section{結論}
本專題以一個輕量、置換等變、成本感知的 LSTM 投資組合網路為核心，透過\emph{注水式溢出重分配}消除結構性現金拖累、\emph{Sharpe 加權集成 + Top-$N$ 截斷}保留高信念部位，並以\emph{擴充候選池 + 強化集中度}進一步提升選股品質。嚴謹的滾動前向驗證顯示 v6（88 檔、Top-10）之 OOS 中位超額達 +5.65\%，明顯優於基線。系統已自動化實盤部署，內建風險停損與最小交易次數機制，完整滿足專題之系統要求。

\appendix
\section{程式架構（節錄）}
\begin{itemize}[leftmargin=2em]
  \item \texttt{stock\_api/}：TWSE／TPEX 抓取、symbol 對映、交易 API 封裝。
  \item \texttt{RL/feature.py}：14 維特徵與橫斷面排名。
  \item \texttt{RL/dl\_portfolio.py}：\texttt{PortfolioNetLSTM}、\texttt{cap\_water\_fill}、\texttt{ensemble\_topn\_truncate}、Sharpe 損失、\texttt{train\_one\_fold\_lstm}。
  \item \texttt{RL/dl\_train\_deploy.py}：5 seed 重訓並輸出自帶設定之檢查點。
  \item \texttt{RL/walk\_forward\_dl\_ensemble.py} / \texttt{probe\_ensemble\_topn.py}：10 折前向驗證與 Top-$N$ 離線掃描。
  \item \texttt{RL/deploy\_rl.py}：對帳、整股下單、回撤停損、湊單（\texttt{run\_pad\_pairs}）、帳本持久化。
  \item \texttt{.github/workflows/}：每日部署與每日重訓之 CI。
\end{itemize}

\end{document}
